% !TEX root = ReadMe.tex
% !TEX TS-program = pdflatex
\documentclass[11pt,a4paper]{article}
\usepackage[utf8]{inputenc}
\usepackage[T1]{fontenc}
\usepackage[ngerman]{babel}
\usepackage{geometry}
\usepackage{listings}
\geometry{margin=2.6cm}
\lstset{basicstyle=\ttfamily\small,breaklines=true,frame=single,columns=fullflexible}
\setlength{\parskip}{0.6em}
\setlength{\parindent}{0pt}
\begin{document}

\section*{JKU Latex template}

Hauptdatei ist \texttt{main.tex}. Dort steht auch \verb|\documentclass|, damit Texifier die Datei als LaTeX erkennt. Diese Anleitung gehört zum Projekt, wird aber nicht in die Arbeit eingebunden. In Texifier liegt sie über \texttt{setup/dependencies.tex} in der Projektliste. Ein Build dieser Datei setzt nur die Anleitung.

\section*{Welche Datei wofür}

\begin{itemize}
\item \texttt{setup/config.tex} -- Deckblatt, Schalter, Buch oder Normal.
\item \texttt{setup/texifier.tex} -- einzeiliger PDF-Titel, Autor, Linkfarben.
\item \texttt{setup/characteristics.tex} -- ausgeschriebene Abschlüsse und Studien. Nur anfassen, wenn ein Wert fehlt.
\item \texttt{subfiles/} -- Kapitel. Jeder Unterabschnitt beginnt mit einem Platzhalter.
\item \texttt{subfiles/general/abstract\_en.tex} und \texttt{abstract\_de.tex} -- Abstract und Kurzfassung.
\item \texttt{subfiles/general/declaration.tex} und \texttt{declaration\_de.tex} -- eidesstattliche Erklärungen.
\item \texttt{references/references.bib} -- Literatur. Im Text mit \verb|\cite{key}|.
\item \texttt{setup/stylesCustom/listings.tex} -- Aussehen und Sprachen der Codeblöcke.
\item \texttt{setup/dependencies.tex} -- Dateien, die Texifier sonst nicht anzeigt.
\end{itemize}

Dateien unter \texttt{setup/styles/} und \texttt{setup/stylesCover/} setzen Seite, Kopfzeile und Deckblatt. Die braucht man für den Text der Arbeit nicht.

\section*{Einstellungen in config.tex}

Schalter sind \verb|\show| oder \verb|\hide|.

\begin{itemize}
\item \verb|\showJKUCover| -- JKU-Deckblatt.
\item \verb|\showAbstract| -- englisches Abstract.
\item \verb|\showAbstractBoth| -- deutsche Kurzfassung. Die Arbeit selbst bleibt Englisch.
\item \verb|\showListOfContent| -- Inhaltsverzeichnis.
\item \verb|\showListOfFigures|, \verb|\showListOfTables|, \verb|\showListOfSourceCode| -- Verzeichnisse im Literaturteil.
\item \verb|\showReferences| -- Literaturverzeichnis.
\item \verb|\showDeclaration| -- beide eidesstattlichen Erklärungen.
\item \verb|\showSubmitter|, \verb|\showSupervisor|, \verb|\showSubmittedAt|, \verb|\showDate| -- die gleichnamigen Zeilen auf dem Deckblatt.
\item \verb|\printLayout| -- \verb|\book| oder \verb|\normal|. Buch: doppelseitig, Kapitel beginnen rechts, leere linke Seite wenn nötig. Normal: einseitig, keine leeren Rückseiten, Kapitel trotzdem auf einer neuen Seite. Die Ränder bleiben gleich.
\item \verb|\type| -- \verb|\DisThesis|, \verb|\DipThesis|, \verb|\MasThesis| oder \verb|\BacThesis|. Steuert die Bezeichnung auf Deckblatt und Erklärung.
\item \verb|\name|, \verb|\supervisor| -- Autor und Betreuer. \verb|\\| bricht eine Zeile auf dem Deckblatt um.
\item \verb|\secondexaminer|, \verb|\assist|, \verb|\company|, \verb|\supervisorComp| -- optional. Leer lassen, dann erscheinen sie nicht.
\item \verb|\degree|, \verb|\study| -- Abschluss und Studium. Fertige Texte stehen in \texttt{characteristics.tex}, zum Beispiel \verb|\DegreeMSc| und \verb|\CompSc|.
\item \verb|\CoverTitle| -- Titel auf dem Deckblatt, mit \verb|\\| umbrochen.
\item \verb|\date|, \verb|\submittedAt| -- Datum und Institut.
\end{itemize}

Der PDF-Titel ist eine einzige Zeile und steht in \texttt{setup/texifier.tex}, zusammen mit \verb|\author| und \verb|\hypersetup|.

\section*{Code}

Einen Block setzt man so. \texttt{language} ist eine der Sprachen unten. \texttt{caption} erscheint im Listingverzeichnis.

\begin{verbatim}
\begin{lstlisting}[language=Kotlin, caption={Beispiel}]
fun main() { println("Hello") }
\end{lstlisting}
\end{verbatim}

Eine ganze Datei statt des Inline-Codes:

\begin{lstlisting}[language=TeX]
\lstinputlisting[language=Java, caption={Orchestrator}]{code/Orchestrator.java}
\end{lstlisting}

Die Datei zusätzlich in \texttt{setup/dependencies.tex} eintragen, sonst zeigt Texifier sie nicht:

\begin{lstlisting}[language=TeX]
%_force-texpad-dependency: code/Orchestrator.java
\end{lstlisting}

Fertige Sprachen, Name in \texttt{language=} so schreiben:

\begin{itemize}
\item C, C++ oder Cpp, CSharp, Java, JavaScript, Python
\item TypeScript, Kotlin, Vue, JSON, Go, Rust
\item HTML, CSS, SQL, PHP, XML, Bash, TeX
\end{itemize}

Eine weitere Sprache kommt an das Ende von \texttt{setup/stylesCustom/listings.tex}. Danach heißt der Block \texttt{language=Swift}. Schlüsselwörter durch Kommas, ohne Leerzeichen am Komma.

\begin{lstlisting}[language=TeX]
\lstdefinelanguage{Swift}{
  morekeywords={class,func,let,var,if,else,return,import,true,false},
  sensitive=true,
  morecomment=[l]{//},
  morecomment=[s]{/*}{*/},
  morestring=[b]"
}
\end{lstlisting}

Schrift, Rahmen und Zeilennummern stehen im \verb|\lstset| derselben Datei und gelten für alle Sprachen.

\section*{Weitere Dateien}

Das Themenblatt liegt als \texttt{Topic/Topic.pdf}. Die Zeile mit \verb|\includepdf| in \texttt{setup/styles/setup\_pages.tex} ist auskommentiert und wird aktiv, sobald die Datei da ist. Der Lebenslauf ist \texttt{CV/cv.pdf}, die Zeile dafür steht in \texttt{setup/styles/setup\_end.tex}.

Ein neues Kapitel ist eine Datei unter \texttt{subfiles/}, begonnen mit \verb|\sectionInit| und \verb|\section{...}|, und eine \verb|\include|-Zeile in \texttt{main.tex}.

\end{document}
